% Fragmento integrable en VII: Relaciones estructurales entre las constantes físicas.
% Insertar después de desarrollos/normalizacion_conjunta.tex.
% Fuentes y cobertura íntegra: editorial/INSERCION_VII.md.
% Sin preámbulo ni macros; usa amsmath y los entornos ya definidos por VII.
\section{Gravité, énergie et autoinertie du système conjoint}
\label{hmtcuatro:gea:seccion}

La normalisation conjointe permet de formuler une question opératorielle : comment la géométrie, la matière et la mémoire interviennent dans une même réalisation. Trois liens sont réunis. Le caractère positif conserve simultanément ses lectures énergétique, massique et radiale ; la variation du transport conjoint produit des courants géométriques et internes d’une même énergie ; et la réponse autoinertielle s’exprime au moyen du rayon gravitationnel de ce caractère. On démontre la conservation de ces liens sous réduction de mémoire et la propagation des contraintes de l’action couplée au cours de son évolution régulière.

Le point de départ est le noyau commun de l’article. APP produit les deux feuillets avec leurs résidus, quotients et incidences ; TRIT conserve régime et orientation : $+1$ correspond au retour et à la mémoire, $0$ au seuil présent et $-1$ à l’ouverture. TPK sélectionne, transporte, plie, retourne et conserve l’état enrichi et son histoire. Le prolongement $w_6\to w_{12}\to w_{18}\to w_{24}\to w_{30}\to R_{36}\to G_9$ maintient préfixe, retenue, survie, route et frontière. La dynamique, son holonomie nonadique et la lecture réduite suivent l’ordre $U_t\to\Gamma_9\to K_{\rm ph}$ : la phase revient, la mémoire avance et l’état n’est pas réinitialisé. Les cinq constructions consubstantielles de la structure discrète du continu ne sont pas séparées lors de la réalisation de leurs lecteurs.

Les constantes et les échelles employées sont des sorties de cette généalogie. La quadrirelation $(\pi,\varphi,e,\alpha)$ conserve son origine commune et son ordre constructif interne ; action, vitesse constitutive et Barbero--Immirzi sont reçus avec leurs lecteurs. Le langage opératoriel et variationnel compose et reconnaît ces sorties, sans introduire de valeurs métrologiques pour sélectionner des états ou des coefficients.

\subsection{Énergie, masse et rayon du même caractère}
\label{hmtcuatro:gea:caracter}

On conserve la réalisation marquée de la section précédente :
\[
\Omega=\{(j,k,q,u,v):j,k\in\mathbb Z/12\mathbb Z,\ 
k-j\in\{0,3,6,9\},\ 1\leq q\leq8,\ 
u<v,\ u,v\in\{1,2,4,5,7,8\}\}.
\]
Son cardinal $N_\Omega=12\cdot4\cdot8\binom62=5760$ compte des incidences, non des états de matière ni des laps. Pour l’inverseur transporté $B:V_{\rm in}\to V_{\rm out}$, avec $B^*B=BB^*=I/4$, le projecteur $P=uu^*$ du mode commun a $|u_i|^2=1/N_\Omega$. L’archive conserve
\[
C_0=(I-P)B,\qquad M_0=PB,\qquad C_0+M_0=B,\qquad\mathsf U=2B.
\]
Les images de $C_0,M_0$ sont orthogonales et $\mathsf U$ est unitaire. L’espérance marquée $\Delta_\Omega(A)=\sum_iP_iAP_i$, $P_i=e_ie_i^*$, est unique parmi les espérances bimodulaires vers l’algèbre diagonale qui fixent les $P_i$ : sur une unité matricielle, la bimodularité exige $\mathcal E(E_{ij})=P_i\mathcal E(E_{ij})P_j$, qui vaut zéro si $i\ne j$ et vaut $E_{ii}$ si $i=j$.

De $C_0C_0^*=(I-P)/4$ et de $P_iPP_i=P_i/N_\Omega$, il résulte
\begin{equation}
\Delta_\Omega(C_0C_0^*)=r_\Omega I,\qquad
\Delta_\Omega(M_0M_0^*)=\frac{I}{4N_\Omega},\qquad
r_\Omega=\frac{5759}{23040}.
\label{hmtcuatro:gea:retorno}
\end{equation}
La somme des deux lectures vaut $I/4$. La règle constitutive longitudinale applique linéairement le premier coefficient à une section de longueur :
\begin{equation}
L=\alpha^{16}r_\Omega L_*,\qquad D=L\mathsf U.
\label{hmtcuatro:gea:longitud}
\end{equation}
La norme locale $\alpha^{16}$ et la section $L_*$ sont des antécédents explicites ; prendre une racine d’intensité constituerait un autre lecteur. La règle métrique radiale est $\|R_Xv\|^2=\|XD^*v\|^2$. Pour $X=X^*>0$ dans la fibre finie, la polarisation et l’unique racine positive donnent $R_X^2=DX^2D^*=L^2\mathsf UX^2\mathsf U^*$ et $R_X=LY$, avec $Y=\mathsf UX\mathsf U^*>0$. Le spectre n’est pas sélectionné à partir d’une réponse gravitationnelle souhaitée.

La même section $\hbar>0$ et l’horloge conjuguée $\omega L=c$ déterminent $E_L=\hbar c/L$ et les lecteurs
\begin{equation}
H_X=E_LY,\qquad M_X=H_X/c^2,\qquad
R_X=LY,\qquad\Lambda_X=LY^{-1}.
\label{hmtcuatro:gea:lectores}
\end{equation}
La notation $H_X$ est réservée à ce lecteur du caractère. L’hamiltonien conjoint $\mathbb H$ introduit ensuite possède un autre domaine et ne s’identifie pas à $H_X$ du seul fait qu’il porte aussi le nom d’énergie.

\begin{proposition}[Coefficient radial commun]
\label{hmtcuatro:gea:coeficiente} Dans l’identification explicite de $R_X$ comme rayon gravitationnel réduit, $R_X=GM_X/c^2$, il existe un unique coefficient compatible :
\begin{equation}
G=\frac{c^3L^2}{\hbar},\qquad
R_X=\frac{G}{c^4}H_X=\frac{G}{c^2}M_X,\qquad
R_X\Lambda_X=L^2I,\quad H_X\Lambda_X=\hbar cI.
\label{hmtcuatro:gea:radial}
\end{equation}
\end{proposition}
\begin{proof}
Multiplier les lecteurs prouve les deux produits. L’identification radiale équivaut à $LY=G\hbar Y/(c^3L)$. Simplifier par le caractère inversible détermine $G$ ; la substitution inverse vérifie l’égalité pour chaque caractère positif. Si deux coefficients la satisfont, leur différence multipliée par $Y$ est nulle et ils coïncident.
\end{proof}

Les règles longitudinale et métrique fixent la réalisation ; la convention du rayon réduit fixe la signification physique de $G$. Le rayon de Schwarzschild de la même masse est $2R_X$. L’égalité conserve les superpositions et les éléments non diagonaux dans toute base transportée, et non les seules valeurs moyennes. Pour une famille différentiable et une connexion de repère commune $\mathcal A_t$, avec $G,c$ fixés,
\begin{equation}
\mathcal D_tR_X=\frac{G}{c^4}\mathcal D_tH_X,\qquad
\mathcal D_tM_X=c^{-2}\mathcal D_tH_X,\qquad
\mathcal D_tO=\dot O+[\mathcal A_t,O].
\label{hmtcuatro:gea:derivada}
\end{equation}
Cela s’obtient en différenciant l’identité complète, repère compris. Cela n’implique pas l’immobilité de l’état.

Pour des modes positifs, ou pour des opérateurs sur des facteurs tensoriels distincts,
\begin{equation}
m(y)=\frac{\hbar}{cL}y,\qquad
\frac{Gm_1m_2}{\hbar c}=y_1y_2,\qquad
\frac{R(y)}{\bar\lambda(y)}=y^2,\quad \bar\lambda(y)=L/y.
\label{hmtcuatro:gea:adimensional}
\end{equation}
La substitution annule les sections dimensionnelles ; les opérateurs sur des facteurs distincts commutent. La taille adimensionnelle dépend du caractère, sans autre paramètre libre par espèce. Il n’en découle ni loi spatiale de force ni évaluation numérique sans évaluer le caractère.

Le calendrier conserve $t_P=L/c$, $t_0=\pi L/(54c)$ et $108t_0=2\pi t_P$. Par conséquent, $L=(54/\pi)\ell_0$, avec $\ell_0=ct_0$ ; les deux longueurs ne sont pas identifiées. Entre les sections $\hbar_b=\rho\hbar_a$, $\rho>0$, avec $L,c,y_i$ fixés, on a $G_b=G_a/\rho$, $m_{i,b}=\rho m_{i,a}$ et la quantité $Gm_1m_2/(\hbar c)$ demeure invariante. On compare des sections conjointes, non une variation temporelle dont les autres coordonnées dimensionnelles seraient artificiellement maintenues fixes.

\subsection{Mémoire statique, dynamique et domaine spectral}
\label{hmtcuatro:gea:memoria}

En dimension finie, ou avec des blocs bornés et coercifs, écrivons
\[
Y=\begin{pmatrix}A&F\\F^*&C\end{pmatrix}>0,\quad
Y_{\rm ef}=A-FC^{-1}F^*,\quad \eta=y+C^{-1}F^*b.
\]
La complétion du carré démontre
\begin{equation}
\left\langle\binom b y,Y\binom b y\right\rangle
=\langle b,Y_{\rm ef}b\rangle+\langle\eta,C\eta\rangle.
\label{hmtcuatro:gea:schur-prueba}
\end{equation}
D’où $Y_{\rm ef}>0$, l’extension minimisante $y=-C^{-1}F^*b$ et la reconstruction $y=\eta-C^{-1}F^*b$. Le résidu $\eta$ est conservé. Comme $\operatorname{Schur}(sY)=sY_{\rm ef}$ pour $s>0$, on obtient $H_{X,\rm ef}=E_LY_{\rm ef}$ et $R_{X,\rm ef}=LY_{\rm ef}$. Résoudre $Y(x,y)=(f,0)$ montre que le bloc de frontière de $Y^{-1}$ est $Y_{\rm ef}^{-1}$, donc $\Lambda_b=LY_{\rm ef}^{-1}$. La compression directe de $Y$ donnerait $A$ : ce n’est pas la même opération.

Soient $\chi=L/E_L=G/c^4>0$ et $R_X=\chi H_X$. Dans une décomposition compatible avec les domaines, nous définissons
\[
\mathcal K_{H_X}(z)=H_{bb}-zI-V(H_{ii}-zI)^{-1}V^*.
\]
\begin{proposition}[Mémoire dynamique complète]
\label{hmtcuatro:gea:resolvente-prop} Dans le domaine résolvant du bloc intérieur,
\begin{equation}
\mathcal K_{R_X}(\chi z)=\chi\mathcal K_{H_X}(z).
\label{hmtcuatro:gea:resolvente}
\end{equation}
Si les deux compléments sont inversibles, leurs inverses portent le facteur $\chi^{-1}$.
\end{proposition}
\begin{proof}
Chaque bloc radial vaut $\chi$ fois le bloc énergétique et $(\chi H_{ii}-\chi zI)^{-1}=\chi^{-1}(H_{ii}-zI)^{-1}$. La substitution laisse un facteur $\chi$ dans chaque terme, en conservant l’ordre. L’inversion prouve la dernière assertion.
\end{proof}
Le paramètre énergétique $z$ est transporté vers la longueur $\chi z$ ; l’identité conserve chaque ordre de mémoire spectrale. Pour $|z|<\|H_{ii}^{-1}\|^{-1}$, la série de la résolvante donne
\[
\mathcal K_{H_X}(z)=H_{\rm est}-zZ+O(z^2),\quad
H_{\rm est}=H_{bb}-VH_{ii}^{-1}V^*,\quad
Z=I+VH_{ii}^{-2}V^*\geq I.
\]
La positivité découle de $Z-I=(H_{ii}^{-1}V^*)^*(H_{ii}^{-1}V^*)$. L’extension statique $(b,-H_{ii}^{-1}V^*b)$ a pour norme au carré $\langle b,Zb\rangle$ : énergie et norme temporelle proviennent de la même extension.

Si $T=Z^{1/2}$ est borné et inversible, le transport est dual :
\begin{equation}
H_{X,c}=T^{-*}H_{X,\rm ef}T^{-1},\quad
R_{X,c}=T^{-*}R_{X,\rm ef}T^{-1},\quad
\Lambda_c=T\Lambda_bT^*.
\label{hmtcuatro:gea:dual}
\end{equation}
La multiplication prouve $R_{X,c}=\chi H_{X,c}$, $R_{X,c}\Lambda_c=L^2I$ et $H_{X,c}\Lambda_c=\hbar cI$, sans commutativité de $T$ avec le caractère. Si $\psi=T(t)b$, $\dot\psi=\dot Tb+T\dot b$ est conservé. Le terme temporel de connexion du générateur n’est pas incorporé rétrospectivement au caractère radial sans déclarer son application.

\begin{proposition}[Extension spectrale et raffinement compatible]
\label{hmtcuatro:gea:limite} Soit $X$ autoadjoint, positif et injectif, et $\mathsf U$ unitaire. Les lecteurs de \eqref{hmtcuatro:gea:lectores}, définis par le calcul spectral de $Y=\mathsf UX\mathsf U^*$, satisfont $R_X=\chi H_X$ sur $\operatorname{Dom}(Y)$. Leurs produits avec $\Lambda_X$ sont valables sur $\operatorname{Dom}(Y^{-1})$ et possèdent les prolongements bornés de \eqref{hmtcuatro:gea:radial}. Si des plongements isométriques satisfont $X'J=JX$, $\mathsf U'J=K\mathsf U$ et conservent $L,E_L$, alors $R'_XK=KR_X$ et $H'_XK=KH_X$ sur les domaines transportés.
\end{proposition}
\begin{proof}
Sur $P_{\varepsilon,M}=\mathbf1_{[\varepsilon,M]}(Y)$, lecteurs et inverses sont bornés et les identités sont des multiplications de fonctions de $Y$. Pour $\psi\in\operatorname{Dom}(Y)$,
\[
\|Y(P_{\varepsilon,M}-I)\psi\|^2
=\int_{(0,\infty)\setminus[\varepsilon,M]}
\lambda^2\,d\langle\psi,E_Y(\lambda)\psi\rangle\longrightarrow0.
\]
L’injectivité élimine la masse spectrale en zéro ainsi que $P_{\varepsilon,M}\psi\to\psi$. Il s’agit d’une convergence en norme de graphe. Si $\psi\in\operatorname{Dom}(Y^{-1})$, alors $Y^{-1}\psi\in\operatorname{Dom}(Y)$ et $YY^{-1}\psi=\psi$, ce qui prouve les produits et leurs prolongements. Substituer les entrelacements dans $R'_X=L\mathsf U'X'\mathsf U'^*$ prouve $R'_XK=KR_X$ ; le calcul énergétique est identique.
\end{proof}
Le secteur nul demeure séparé. On précise ainsi l’extension non bornée de la normalisation précédente. Une projection orthogonale arbitraire n’est pas pour autant admise dans les formules par blocs : pour les opérateurs non bornés, des domaines compatibles ou des formes fermées à blocs contrôlés sont requis.

\subsection{Rétroaction dans le même transport}
\label{hmtcuatro:gea:retroaccion}

Le registre géométrique $x$, les liens internes $u$ et les champs $f$ proviennent du même état enrichi. Dans une cellularisation orientée,
\begin{equation}
E(f,x,u)=\langle r,W(x,u)r\rangle,\quad r=D_Tf,\quad
(D_Tf)_a=f_{t(a)}-T_af_{s(a)},\quad
T_a=U_a^{\rm sp}(x)\otimes R_{\rm int}(u_a).
\label{hmtcuatro:gea:energia-memoria}
\end{equation}
$W=W^*>0$ conserve les blocs entre incidences. Spin et représentation interne possèdent des facteurs distincts ; la réalisation chirale conserve ses blocs et l’application Higgs--Yukawa équivariante.

\begin{proposition}[Différentielle d’une même énergie]
\label{hmtcuatro:gea:variacion} À champs fixés et avec $q=Wr$,
\begin{align}
\delta E&=-2\operatorname{Re}\sum_a
\langle q_a,\delta T_af_{s(a)}\rangle+\langle r,\delta Wr\rangle,
\label{hmtcuatro:gea:diferencial}\\
\delta T_a&=\delta U_a^{\rm sp}\otimes R_{\rm int}(u_a)
+U_a^{\rm sp}\otimes\delta R_{\rm int}(u_a).\nonumber
\end{align}
\end{proposition}
\begin{proof}
Dériver la forme quadratique donne deux termes conjugués par $W=W^*$ et le terme $\langle r,\delta Wr\rangle$. Utiliser $\delta r_a=-\delta T_af_{s(a)}$ prouve la première formule ; la règle du produit prouve la seconde. Les deux variations agissent sur les mêmes $r,q$, sans courants choisis indépendamment.
\end{proof}

\subsubsection{Un espace quantique et un opérateur conjoint}
\label{hmtcuatro:gea:cuantica}

On fixe une cellularisation finie, un ensemble fini de registres géométriques $\mathcal X$ et un espace de Fock fermionique fini $\mathcal F$. Les liens parcourent le produit compact $\mathcal G^E$ reçu ; les doublets $\Phi_v\in\mathbb C^2$ n’ont pas de coupure d’amplitude. Les mesures invariantes $d\mu_x=Z_x^{-1}e^{-S_x}d\mu_0$ possèdent des densités positives bornées à inverse borné et ne dépendent pas de $\Phi$ dans cette carte. L’espace est
\[
\mathscr H=\bigoplus_{x\in\mathcal X}
L^2(\mathcal G^E\times\mathbb C^{2|V|},
\mu_x\otimes d\Phi;\mathcal F).
\]
Les identifications $A_xf=\sqrt{Z_x}e^{S_x/2}f$ sont unitaires à partir de la mesure commune et transportent l’horloge vers $H_{\rm clk}^{\mu}=A(H_{\rm clk}\otimes I)A^*$. L’horloge mélange des registres géométriques du même espace. L’opérateur est réalisé au moyen de la forme de
\begin{equation}
\mathbb H=H_{\rm clk}^{\mu}+\hbar cK_{\rm gauge}
+d\Gamma(D_T^*WD_T)+H_{\Phi,\rm cin}+V_\Phi
+\widehat H_Y+\widehat H_{\rm tor}.
\label{hmtcuatro:gea:total}
\end{equation}
L’application de Yukawa est linéaire en $\Phi,\bar\Phi$ et satisfait $Y(g\Phi)\rho_R(g)=\rho_L(g)Y(\Phi)$. La contorsion éliminée n’est pas conservée simultanément comme variable indépendante dans les termes de base de cette carte effective.

La cinétique des moments scalaires a la forme $\int\langle\nabla_\Phi\Psi,A(x,u)\nabla_\Phi\Psi\rangle$, avec $A$ bornée, uniformément positive et indépendante de $\Phi$. On exige $A(x,u^g)=O_gA(x,u)O_g^{\mathsf T}$ pour l’action réelle $O_g$ des doublets ; la positivité seule ne prouve pas cette covariance. L’énergie spatiale scalaire est un potentiel quadratique positif majoré par $C_{\rm grad}r^2$, $r^2=\sum_v|\Phi_v|^2$, non un opérateur borné sur tout $L^2$. Les volumes satisfont $0<w_{\min}\leq w_v(x)\leq w_{\max}$ et $\lambda_\Phi>0$. L’horloge, le bloc de jauge, les transports, $W$ et les coefficients torsionnels sont bornés dans la carte. Le potentiel conserve sa valeur absolue :
\begin{equation}
V_\Phi=\sum_vw_v\lambda_\Phi(|\Phi_v|^2-v_0^2/2)^2
-\sum_vw_v\lambda_\Phi v_0^4/4.
\label{hmtcuatro:gea:potencial}
\end{equation}
La dernière contribution dépend du volume et n’est pas soustraite.

\begin{theorem}[Forme fermée et évolution conjointe]
\label{hmtcuatro:gea:forma} La forme de \eqref{hmtcuatro:gea:total} est fermée et minorée sur $\mathcal Q=\{\Psi\in\mathscr H:\nabla_\Phi\Psi\in L^2,\
r^2\Psi\in L^2\}$. Elle détermine un unique opérateur autoadjoint associé $\mathbb H$ et l’évolution unitaire $e^{-it\mathbb H/\hbar}$.
\end{theorem}
\begin{proof}
Cauchy--Schwarz donne $\sum_vw_v\lambda_\Phi|\Phi_v|^4\geq ar^4$, $a=\lambda_\Phi w_{\min}/|V|>0$. L’espace de Fock fini et Yukawa linéaire donnent $\|\widehat H_Y(\Phi)\|\leq C_Yr$ ; la torsion est bornée. Pour $\varepsilon>0$,
\[
C_Yr\leq\varepsilon r^4+
\frac{3C_Y^{4/3}}{4^{4/3}\varepsilon^{1/3}},\qquad
br^2\leq\varepsilon r^4+\frac{b^2}{4\varepsilon}.
\]
La première borne maximise $C_Yr-\varepsilon r^4$ ; la seconde complète le carré en $r^2$. Les appliquer avec $\varepsilon=a/4$ à Yukawa et au terme quadratique négatif donne
\[
q(\Psi)\geq c_1\|\nabla_\Phi\Psi\|^2+
(a/2)\|r^2\Psi\|^2-C\|\Psi\|^2,\qquad c_1>0.
\]
La norme $\|\Psi\|^2+\|\nabla_\Phi\Psi\|^2+\|r^2\Psi\|^2$ est complète par fermeture de la dérivée faible et de la multiplication. Ces bornes et les majorations rendent la norme de forme, après ajout d’une constante, équivalente à cette norme complète. Le théorème de représentation des formes fermées produit l’opérateur associé et son calcul spectral donne l’évolution. L’unicité concerne cette forme et ce domaine, non chaque domaine formel alternatif.
\end{proof}

La covariance de $D_T,W$, l’invariance des mesures, l’identité de Yukawa et la contraction torsionnelle interne conservent $q$ et $\mathcal Q$. L’unicité de l’opérateur associé implique que la représentation compacte commute avec $\mathbb H$ ; sa moyenne de Haar réduit l’opérateur et son évolution. Ce n’est pas une moyenne corrective ultérieure.

On retire la coupure computationnelle au moyen de blocs complets de Peter--Weyl et de couches isotropes complètes d’Hermite. La régularisation et la troncature lisse approchent les dérivées et le poids $r^2$ ; les densités bornées préservent des normes équivalentes. Leur réunion est dense dans $\mathcal Q$. Pour les projections orthogonales $P_n$ dans les mesures $\mu_x$ et les opérateurs de Galerkin $\mathbb H_n$,
\[
(\mathbb H_n+\lambda)^{-1}P_nf\longrightarrow
(\mathbb H+\lambda)^{-1}f\qquad(\lambda>C).
\]
L’équation coercive de forme et sa version de Galerkin donnent l’orthogonalité de l’erreur ; la continuité et la coercivité la majorent par la meilleure erreur d’approximation, qui tend vers zéro. On retire les coupures scalaire et fonctionnelle des liens de cette carte, non les coupures spatiale, géométrique ou fermionique. L’autoadjonction n’implique pas non plus que $e^{-t\mathbb H}$ soit à trace.

\subsubsection{Échange et conservation}
\label{hmtcuatro:gea:intercambio}

Pour $V=\sum_xP_x\otimes B_x$ et un élément $h_{xy}$ de l’horloge, la multiplication par blocs donne
\begin{equation}
[H_{\rm clk}\otimes I,V]_{xy}=h_{xy}(B_y-B_x).
\label{hmtcuatro:gea:conmutador}
\end{equation}
Si $h_{xy}\ne0$ et $B_y\ne B_x$, il y a échange géométrique--matériel. Si $\mathbb H=A+B+V$, pour des opérateurs bornés ou un cœur commun justifiant les dérivées,
\[
\frac d{dt}\langle A\rangle=\frac i\hbar\langle[\mathbb H,A]\rangle,
\quad
\frac d{dt}\langle B\rangle=\frac i\hbar\langle[\mathbb H,B]\rangle,
\quad
\frac d{dt}\langle V\rangle=\frac i\hbar\langle[\mathbb H,V]\rangle.
\]
La somme est nulle par $[\mathbb H,\mathbb H]=0$. L’énergie est conservée, interaction comprise, non chaque secteur séparément. La géométrie est dans l’espace quantique, non ajoutée comme fond extérieur. Les histoires étendues conservent mémoire et retour ; leur raffinement n’équivaut pas par définition au raffinement spatial.

\subsection{La même action et ses courants}
\label{hmtcuatro:gea:accion}

La réalisation géométrique emploie une cotétrade non dégénérée $e$, une connexion de Lorentz $\omega$ et des connexions internes de la même matière. Soient $\Sigma^{IJ}=e^I\wedge e^J$, $P_\gamma=\star+\gamma^{-1}I$, $\gamma\in\mathbb R\setminus\{0\}$ et $\star^2=-I$ sur les bivecteurs. Les couplages restent fixes pendant la variation. Avec $\kappa=8\pi G/c^4$, la même action est
\begin{equation}
S_{\rm tot}=\frac{\hbar}{16\pi L^2}
\int\Sigma_{IJ}\wedge(P_\gamma F_\omega)^{IJ}
-\frac{\Lambda\hbar}{8\pi L^2}\int\operatorname{vol}_e+S_m,
\label{hmtcuatro:gea:accion-total}
\end{equation}
où $S_m=S_{\rm YM}+S_\Phi+S_{\rm Weyl}+S_Y$ contracte ses indices avec la même $e$, que l’on fait varier. Les coefficients découlent de $(2\kappa c)^{-1}=\hbar/(16\pi L^2)$.

Avec $\delta_\omega S_m=-\tfrac12\int
\delta\omega^{IJ}\wedge\sigma_{IJ}$, $\delta F=D_\omega\delta\omega$ et l’intégration par parties,
\begin{equation}
D_\omega(P_\gamma\Sigma)=\kappa c\sigma.
\label{hmtcuatro:gea:cartan}
\end{equation}
On utilise des variations intérieures ou la complétion de bord compatible. Le courant de spin et la source métrique sont des dérivées de la même matière, non deux entrées indépendantes.

\subsubsection{Élimination torsionnelle et termes croisés}
\label{hmtcuatro:gea:torsion}

L’inverse réel est $P_\gamma^{-1}=\gamma^2(-\star+\gamma^{-1}I)/(1+\gamma^2)$ : le produit des deux polynômes en $\star$ avant normalisation vaut $(1+\gamma^{-2})I$. Notons $\mathcal Y^{IJ}=\kappa c(P_\gamma^{-1}\sigma)^{IJ}$ la forme à valeurs bivectorielles de degré trois, distincte du caractère $Y$. Pour $T^I=D_\omega e^I$, $D_\omega\Sigma^{IJ}=T^I\wedge e^J-e^I\wedge T^J$. Avec le repère dual $E_I$, son inverse est
\[
\mathcal B^I=\sum_J\iota_{E_J}\mathcal Y^{IJ},\qquad
\vartheta=\tfrac14\sum_I\iota_{E_I}\mathcal B^I,\qquad
T^I=\mathcal B^I-e^I\wedge\vartheta.
\]
En effet, si $t=\sum_J\iota_{E_J}T^J$, les contractions donnent $\mathcal B^I=T^I+e^I\wedge t$ puis $4t$. L’application a un noyau nul entre espaces de dimension vingt-quatre ; c’est donc un isomorphisme. Pour $\omega=\mathring\omega(e)+\mathfrak k$,
\[
T_{IJK}=\mathfrak k_{IKJ}-\mathfrak k_{IJK},\qquad
\mathfrak k_{IJK}=\tfrac12(T_{KIJ}-T_{IJK}-T_{JKI}).
\]
La combinaison cyclique et l’antisymétrie dans les deux premiers indices prouvent ces formules. Pour une matière affine en $\mathfrak k$, avec un courant qui en est indépendant, il existe une unique $\mathfrak k_*[e,\sigma]$ linéaire en le courant total $\sigma=\sum_s\sigma_s$.

Le développement de la courbure conserve le bord
\[
S_\partial[e,\mathfrak k]=\frac1{2\kappa c}
\int_{\partial M}\Sigma_{IJ}\wedge(P_\gamma\mathfrak k)^{IJ}.
\]
La densité intérieure dépendant de la contorsion est $\mathcal Q_{e,\gamma}(\mathfrak k)
-\tfrac12\mathfrak k^{IJ}\wedge\sigma_{IJ}$, où $\mathcal Q$ est homogène de degré deux. Sa stationnarité évaluée en $\delta\mathfrak k=\mathfrak k_*$ donne $2\mathcal Q(\mathfrak k_*)=
\tfrac12\mathfrak k_*^{IJ}\wedge\sigma_{IJ}$. Par conséquent,
\begin{equation}
\mathcal L_{\rm spin}^{\rm ef}
=-\tfrac14\mathfrak k_*^{IJ}\wedge\sigma_{IJ}
=-\mathcal Q_{e,\gamma}(\mathfrak k_*).
\label{hmtcuatro:gea:spin}
\end{equation}
La linéarité conserve les termes croisés entre espèces. La contorsion est éliminée une fois ; on n’ajoute pas simultanément son interaction linéaire éliminée et le terme quartique effectif.

Dans la spécialisation spinorielle, le coefficient reçu $C_\gamma=3\kappa c\hbar^2\gamma^2/[16(1+\gamma^2)]$ donne
\begin{equation}
\frac{C_\gamma}{\hbar}
=\frac{3\pi}{2}L^2\frac{\gamma^2}{1+\gamma^2}.
\label{hmtcuatro:gea:coeficiente-torsion}
\end{equation}
Il suffit de substituer $\kappa=8\pi G/c^4$ et $\hbar G=c^3L^2$, sans changer la branche de $\gamma$. Dans l’espace de Fock de la carte, le courant total conserve la contraction
\[
\widehat Q_{{\rm tot},v}=\sum_{a=1}^4s_a
\{d\Gamma(\widetilde M_{a,v})^2-d\Gamma(\widetilde M_{a,v}^2)\},
\quad s=(-1,-1,-1,+1),
\]
\[
\widehat H_{\rm tor}
=-\sum_v\frac{N_vcC_\gamma}{\mathcal V_v}\widehat Q_{{\rm tot},v}.
\]
Les matrices $\widetilde M_{a,v}$ représentent les composantes du courant. $N_v$ est le lapse, distinct de $N_\Omega$ et d’une occupation. L’identité d’ordre normal retire la contraction à une particule, non les termes croisés entre espèces. La signature demeure dans $s_a$.

\subsubsection{Source métrique et transformation de Legendre de la même action}
\label{hmtcuatro:gea:legendre}

L’élimination stationnaire préserve la variation des autres variables : la règle de la chaîne contient un terme en $\delta\mathfrak k_*$ multiplié par l’équation de connexion, qui est nulle. L’action réduite est
\[
S_{\rm red}=\frac1{2\kappa c}\int(R[g]-2\Lambda)\operatorname{vol}_e
+S_m^{\rm LC}+S_{\rm spin}^{\rm ef}+S_\partial.
\]
Nous définissons la source d’action par $\delta_g(S_m^{\rm LC}+S_{\rm spin}^{\rm ef})
=-\tfrac12\int\mathcal T^S_{\mu\nu}\delta g^{\mu\nu}
\operatorname{vol}_e$. La variation gravitationnelle intérieure produit
\begin{equation}
G_{\mu\nu}+\Lambda g_{\mu\nu}
=\kappa c\mathcal T^S_{\mu\nu}
=\kappa(T^{\rm LC}_{\mu\nu}+U^{\rm spin}_{\mu\nu}),
\quad T^{\rm LC}+U^{\rm spin}=c\mathcal T^S.
\label{hmtcuatro:gea:metrica}
\end{equation}
La variation de volume du terme constant du potentiel est incluse. Le facteur $c$ correspond à la cotétrade dans la droite de longueur ; les sources d’action et les sources physiques ne sont pas mélangées.

Avec $a=(2\kappa c)^{-1}$ et $\mathcal B_{IJ}=(P_\gamma\Sigma)_{IJ}$, le potentiel géométrique est $\theta_{\rm geom}=a\mathcal B_{IJ}\wedge\delta\omega^{IJ}$. Décomposer $\omega=\omega_0dt+\underline\omega$ et $F_{0a}=\dot\omega_a-D_a\omega_0$ identifie le terme cinétique ; intégrer par parties conserve Gauss et le bord. Avec $e_0=Nn+N^ae_a$ et la transformation de Legendre de la même matière,
\begin{equation}
S_{\rm tot}=\int dt\,[\Theta_{\rm red}(\dot z)
-\mathcal C[N]-\mathcal V[N^a]-\mathcal G_L[\omega_0]
-\mathcal G_{\rm int}[A_0]]+S_\partial.
\label{hmtcuatro:gea:canonica}
\end{equation}
Lapse, shift et jauge conservent leurs directions dégénérées. Les contraintes de seconde classe sont réduites avant d’utiliser $\Theta_{\rm red}$ ou sont expressément conservées. Les spineurs du premier ordre ne reçoivent pas d’inverse fictif des vitesses.

Pour la contorsion et son moment, les contraintes sont $\pi_{\mathfrak k}=0$ et $\partial\mathcal Q/\partial\mathfrak k-\sigma/2=0$. Leur matrice possède les blocs $\left(\begin{smallmatrix}0&-M\\M^{\mathsf T}&B\end{smallmatrix}\right)$, avec $M$ inversible par la reconstruction précédente ; elle est donc inversible. Le crochet de Dirac effectue cette élimination et conserve séparément les contraintes fermioniques graduées. Dans une direction bosonique régulière,
\[
H_\Phi=N\left(\frac{p^2}{2w}+wV+H_{\rm grad}\right),\quad
p=\frac wN D_t\phi,\quad
L_\Phi=\frac{w}{2N}|D_t\phi|^2-NwV-NH_{\rm grad}.
\]
Le volume $w=\sqrt{\det q}$ n’est pas figé. Pour $K_{ij}=(\dot q_{ij}-\mathcal L_{N^a}q_{ij})/(2N)$,
\[
\Pi^{ij}=\frac{\sqrt q}{2\kappa c}(K^{ij}-q^{ij}K),\qquad
K_{ij}=\frac{2\kappa c}{\sqrt q}
(\Pi_{ij}-\tfrac12q_{ij}\Pi).
\]
La trace et la partie sans trace prouvent l’inverse. Si le moment total contient une contribution matérielle de repère, on utilise $\Pi=p-p_m$, en conservant $p_m$. On retrouve ainsi la même action et sa source. La transformation de Legendre ne démontre pas à elle seule qu’une quantification arbitraire de \eqref{hmtcuatro:gea:canonica} coïncide avec \eqref{hmtcuatro:gea:total}.

\subsection{Propagation des contraintes avec la source totale}
\label{hmtcuatro:gea:restricciones}

Soit $E_{\mu\nu}=G_{\mu\nu}+\Lambda g_{\mu\nu}
-\kappa c\mathcal T^S_{\mu\nu}$. Sur les équations matérielles, l’invariance par difféomorphismes de la fonctionnelle réduite implique $\nabla^\mu\mathcal T^S_{\mu\nu}=0$ : une variation intérieure engendrée par $\xi$ laisse l’intégrale de $\mathcal T^S_{\mu\nu}\nabla^\mu\xi^\nu$, tandis que les équations matérielles annulent les autres termes. Intégrer par parties et faire varier $\xi$ prouve la divergence nulle, avant d’imposer $E=0$. Bianchi et les couplages fixes donnent $\nabla_\mu E^{\mu\nu}=0$.

\begin{proposition}[Conservation des données de contrainte]
\label{hmtcuatro:gea:propagacion} Une solution régulière des équations matérielles et d’évolution métrique conserve les contraintes satisfaites initialement pendant son intervalle d’existence régulière.
\end{proposition}
\begin{proof}
En coordonnées gaussiennes $g=-d\tau^2+q_{ij}(\tau,x)dx^idx^j$, on impose $E_{ij}=0$ et l’on écrit $C=E_{00}$, $C_i=E_{0i}$, $K_{ij}=\dot q_{ij}/2$ et $\mathcal K=q^{ij}K_{ij}$. On a $E^{00}=C$, $E^{0i}=-C^i$, $E^{ij}=0$, $\Gamma^i{}_{0j}=K^i{}_j$ et $\Gamma^\mu{}_{\mu0}=\mathcal K$. Les composantes temporelle et spatiale de la divergence donnent
\[
\partial_\tau C-D_iC^i+\mathcal KC=0,\qquad
\partial_\tau C^i+\mathcal KC^i+2K^i{}_jC^j=0.
\]
En abaissant l’indice, $\dot q_{ij}=2K_{ij}$ annule le dernier terme spatial. Il reste exactement
\begin{equation}
\partial_\tau C-D_iC^i+\mathcal KC=0,\qquad
\partial_\tau C_i+\mathcal KC_i=0.
\label{hmtcuatro:gea:propagacion-sistema}
\end{equation}
La seconde équation homogène maintient $C_i=0$ et la première devient $\dot C+\mathcal KC=0$, en conservant $C=0$. $q$ et $\mathcal K$ sont ceux de la solution couplée, non un fond prescrit.
\end{proof}
La carte gaussienne prouve localement une assertion tensorielle, sans figer $q^{-1}$ ni remplacer la source. On n’affirme pas d’existence globale sans singularités pour chaque donnée. La propagation variationnelle n’est pas non plus à elle seule une identité de commutateurs quantiques pour des déformations normales arbitraires.

\subsection{Frontière et opérateur autoinertiel}
\label{hmtcuatro:gea:autoinercia}

Soient $b:\mathcal X_\infty^{\rm enr}\to B_{\rm ef}$, $B_{\rm ef}=\operatorname{im}(b)$. Une réponse $I_v$ se factorise de manière unique sous la forme $I_v=\overline I_v\circ b$ si et seulement si elle est constante sur les fibres de $b$. La nécessité est immédiate ; pour la suffisance, $\overline I_v(\beta)=I_v(x)$ avec $b(x)=\beta$ est indépendant du représentant. L’image effective donne l’unicité. La dépendance machienne forte ajoute des états ayant la même restriction locale et des réponses distinctes.

La charnière du même état produit $r_{\rm av}=\pi/\varphi^2$ et $r_{\rm av}^J=\varphi/\pi^2$. Dans le secteur d’amplitude commune, $\mathcal L_{\rm ar}=\pi\mu_{\rm ar}\mathcal M(x)$ et $\mathcal L_{\rm vol}=\mu_{\rm vol}\mathcal M(x)$, $\mathcal M(x)>0$, avec $\mu_{\rm ar}/\mu_{\rm vol}=\varphi^{-2}$. Diviser par $\mathcal L_{\rm vol}$ donne des lecteurs du même type $A_\partial=r_{\rm av}$ et $V_\partial=1$. Le quotient annule l’amplitude, non l’histoire conservée dans l’archive.

Pour une orientation $a$, $q_a(x)=h_a(x)/(108t_0(x))=\hbar_a(x)/t_P(x)$ descend à $Q:B_{\rm ef}\to\mathcal L_E^+$ si et seulement si $b(x)=b(x')$ implique $h_a(x)/t_0(x)=h_a(x')/t_0(x')$. C’est le critère de factorisation appliqué à l’horloge. Les cartes fixes et les cartes variables dont le quotient descend le satisfont. Dans la carte radiale, $Q=\hbar_ac/L$. Avec des énergies de cette même droite, $E_1E_2/Q^2$ est adimensionnel et indépendant de la base.

Pour des projections orthogonales complémentaires,
\begin{align}
w_A&=\frac{A_\partial r_{\rm av}}
{A_\partial r_{\rm av}+V_\partial r_{\rm av}^J}
=\frac{\pi^4}{\pi^4+\varphi^5},\qquad w_V=1-w_A,
\label{hmtcuatro:gea:pesos}\\
W_{\rm av}&=w_AP_A+w_VP_V,\qquad
\overline I_v=\frac{E_1E_2}{Q^2}W_{\rm av}.
\label{hmtcuatro:gea:respuesta}
\end{align}
Multiplier le numérateur et le dénominateur de $r_{\rm av}^2/(r_{\rm av}^2+r_{\rm av}^J)$ par $\varphi^4\pi^2$ prouve l’évaluation. La distinction entre lecteur aréal et pondération est conservée : la marque régionale est générée une fois. $0<w_A,w_V<1$ rend $W_{\rm av}$ positif et inversible, et
\[
\langle z,\overline I_vz\rangle
=\frac{E_1E_2}{Q^2}
(w_A\|P_Az\|^2+w_V\|P_Vz\|^2)>0\quad(z\ne0).
\]
Échanger complètement $(A_\partial,r_{\rm av},P_A)$ et $(V_\partial,r_{\rm av}^J,P_V)$ permute les termes. Un transport unitaire conjugue la réponse et conserve son spectre ; réaliser l’échange au sein d’un espace exige des dimensions sectorielles égales. Une remise à l’échelle commune de $A_\partial,V_\partial$ ne change pas les poids et ne prouve pas de dépendance globale. À énergies et poids fixés, $Q\mapsto qQ$ modifie la réponse de $q^{-2}$ ; le témoin fort conserve en outre la même restriction locale.

\begin{theorem}[Autoinertie et rayon du même caractère]
\label{hmtcuatro:gea:autoinercial-radial} Avec la même horloge et la même orientation,
\begin{equation}
\overline I_v=\frac{G_aE_1E_2}{\hbar_ac^5}W_{\rm av}
=\frac{G_am_1m_2}{\hbar_ac}W_{\rm av}.
\label{hmtcuatro:gea:mach-gravedad}
\end{equation}
Pour une sonde $y_p>0$, $E_p=Qy_p$ et $\bar\lambda_p=L/y_p$, l’extension spectrale satisfait
\begin{equation}
\widehat I_{p,Y}=y_pY\otimes W_{\rm av}
=\frac{E_p}{Q^2}H_X\otimes W_{\rm av}
=\frac{R_X}{\bar\lambda_p}\otimes W_{\rm av}
=\frac{G_am_p}{\hbar_ac}M_X\otimes W_{\rm av}.
\label{hmtcuatro:gea:identidad-autoinercial}
\end{equation}
\end{theorem}
\begin{proof}
$Q^2=\hbar_a^2c^2/L^2$ et $\hbar_aG_a=c^3L^2$ donnent $Q^{-2}=G_a/(\hbar_ac^5)$ ; on utilise ensuite $E_i=m_ic^2$. Pour $Y=\sum_jy_jP_j$, appliquer le lecteur à chaque énergie $Qy_j$ produit $\sum_jy_py_jP_j\otimes W_{\rm av}$, indépendamment de la base propre. Le calcul spectral étend le lecteur linéaire au caractère autoadjoint positif général. La décomposition par $P_A,P_V$ donne les opérateurs $y_pw_AY$ et $y_pw_VY$, avec leurs multiplicités : ils sont autoadjoints positifs sur le domaine de $Y\otimes I$. Substituer les lecteurs radiaux prouve les autres expressions.
\end{proof}

Si une réalisation supplémentaire identifie effectivement son opérateur complet à $H_X=QY$, elle conserve cette identification et ses domaines. Pour $H_X=\sum_rH_r$ sur un cœur commun, $\widehat I_{p,Y}=(E_p/Q^2)\sum_rH_r\otimes W_{\rm av}$. La positivité de chaque terme n’est pas requise ; celle du caractère de l’opérateur identifié l’est. Aucun terme n’est supprimé et l’origine de l’énergie n’est pas modifiée pour l’imposer. Le théorème n’attribue pas automatiquement ce lecteur à $\mathbb H$.

Dans les blocs de \eqref{hmtcuatro:gea:schur-prueba}, l’inverse intérieur de $y_pY\otimes W_{\rm av}$ est $y_p^{-1}C^{-1}\otimes W_{\rm av}^{-1}$. La multiplication donne
\begin{equation}
\operatorname{Schur}(\widehat I_{p,Y})
=y_pY_{\rm ef}\otimes W_{\rm av}
=\frac{R_{X,\rm ef}}{\bar\lambda_p}\otimes W_{\rm av}.
\label{hmtcuatro:gea:schur-autoinercia}
\end{equation}
Les blocs croisés ne sont pas effacés. Pour $a_s=y_pw_s>0$, $s=A,V$, le même calcul donne $\mathcal K_{a_sY}(a_sz)=a_s\mathcal K_Y(z)$ dans le domaine résolvant intérieur, en conservant la mémoire complète. Si $Y'K=KY$ et $W'_{\rm av}J=JW_{\rm av}$, avec $y_p$ et la section intensive fixés,
\begin{equation}
\widehat I'_{p,Y'}(K\otimes J)
=(K\otimes J)\widehat I_{p,Y}.
\label{hmtcuatro:gea:naturalidad}
\end{equation}
Cela se prouve par substitution. Les coupures spectrales de la proposition \ref{hmtcuatro:gea:limite} convergent en norme de graphe ; les deux poids fixes transmettent la convergence à la réponse. La naturalité est celle de ces entrelaceurs effectifs, non celle d’une famille distincte d’applications spatiales identifiée par son seul nom.

\subsection{Autoinertie de l’état et accélération de sa projection}
\label{hmtcuatro:gea:proyeccion}

Dans une réalisation différentiable, l’état conjoint $\Gamma=(x,\mathcal O,m)$ possède la connexion $\nabla^{\rm tot}$ sur $\mathcal N$. Une projection lisse $\pi_S:\mathcal N\to\mathcal N_S$, avec la connexion $\nabla^S$, conserve les transports et les lecteurs qu’elle réalise. Le repère observé est $\mathfrak F_{\mathcal O}(\Gamma)
=\operatorname{Res}_{\mathcal O}(\Gamma,\nabla^{\rm tot})$. Pour une courbe deux fois différentiable et $\gamma_S=\pi_S\circ\Gamma$, l’autoinertie signifie $\nabla^{\rm tot}_{\dot\Gamma}\dot\Gamma=0$. Le défaut est
\[
\mathcal K_{\pi_S}(\dot\Gamma,\dot\Gamma)
=\nabla^S_{\dot\gamma_S}\dot\gamma_S
-d\pi_S(\nabla^{\rm tot}_{\dot\Gamma}\dot\Gamma).
\]
\begin{proposition}[Accélération issue de la lecture]
\label{hmtcuatro:gea:aceleracion}
Pour une trajectoire auto-inertielle,
\begin{equation}
a_S=\mathcal K_{\pi_S}(\dot\Gamma,\dot\Gamma),\qquad
(\mathcal K_{\pi_S})^a{}_{BC}
=\partial_B\partial_C\pi_S^a
+(\Gamma^S)^a{}_{bc}\partial_B\pi_S^b\partial_C\pi_S^c
-\partial_D\pi_S^a(\Gamma^{\rm tot})^D{}_{BC}.
\label{hmtcuatro:gea:defecto}
\end{equation}
\end{proposition}
\begin{proof}
La règle de la chaîne donne $\ddot\gamma_S^a=\partial_B\pi_S^a\ddot\Gamma^B
+\partial_B\partial_C\pi_S^a\dot\Gamma^B\dot\Gamma^C$. Ajouter la connexion du sous-système et soustraire l’image de l’accélération totale annule les termes en $\ddot\Gamma$. Les termes restants sont le tenseur affiché contracté avec les vitesses. L’autoinertie annule l’accélération totale et laisse l’égalité.
\end{proof}
La réponse géométrique se factorise par $b$ lorsque la connexion, la projection et la donnée cinématique descendent à la frontière, ou lorsque cette dernière est fixée dans la comparaison. Si la vitesse change, elle est conservée comme argument. L’ensemble autoinertiel peut ainsi donner en lecture des sous-systèmes accélérés sans repère absolu extérieur.

L’autoinertie n’annule pas à elle seule la courbure ni l’holonomie. Une donnée transportée le long d’une boucle revient si et seulement si elle appartient à l’espace fixe de son holonomie. Le retour de l’état complet exige en outre une réalisation injective dans le secteur comparé : retour visible et retour de l’histoire demeurent distincts.

\subsection{Résultat et contrôles de portée}
\label{hmtcuatro:gea:conclusion}

Le même caractère fixe masse, énergie et rayon ; leur proportionnalité conserve mémoire statique, mémoire spectrale et normalisation duale. Les transports géométrique et interne interviennent dans la différentielle d’une seule énergie. L’action conjointe produit les sources géométrique et torsionnelle et conserve les contraintes sur sa solution régulière. L’opérateur autoinertiel est le rayon du caractère divisé par la longueur conjuguée de la sonde, tandis que l’accélération est le défaut de projection de l’état.

Les deux derniers objets ont des types distincts : $\widehat I_{p,Y}$ ne s’identifie ni à $\mathcal K_{\pi_S}$ ni au tenseur d’Einstein. L’incorporation démontrée n’est pas non plus remplacée par l’exigence différente d’annuler une courbure quantique totale pour des déformations normales arbitraires. Cet énoncé possède ses générateurs et ses domaines ; il n’est pas affirmé ici, ni déclaré absent du corpus au motif qu’il ne constitue pas cette conclusion.

Les falsificateurs sont concrets. Changer seulement $G$ ou $\hbar$ avec $L,Y$ fixés rompt \eqref{hmtcuatro:gea:radial} ; remplacer $L$ par $\ell_0$ sans $54/\pi$ altère l’horloge ; effacer $F$ modifie la réduction ; remettre conjointement à l’échelle les lecteurs aréal et volumétrique ne prouve pas la séparation machienne ; une application qui n’entrelace pas $Y$ ne satisfait pas \eqref{hmtcuatro:gea:naturalidad}. Les essais rationnels sont des contrôles finis reproductibles, non des valeurs cibles du générateur ni des substituts aux preuves de domaine et de limite.
